%------------------------------------------------------------------------------
\subsection{Part III: Spectral Differentiation as Biomarker}
\label{sec:biomarker}
%------------------------------------------------------------------------------

A biomarker must first be reliable before it can be meaningful. We therefore begin with temporal stability, then establish the developmental trajectory---the strongest and most replicated signal---before turning to cognitive associations, psychopathology, and informative nulls. The enrichment features derived from logarithmic $\phisym$-lattice coordinates carry systematic variance across lifespan development, cognitive performance, and psychopathology dimensions while remaining silent for personality, handedness, and metabolic variables.

\textbf{Gamma caveat.} Throughout this section, gamma-band results (32--52~Hz) should be interpreted with the scalp EEG contamination caveat described in Section~\ref{sec:two-regimes}: gamma had the weakest cross-dataset consistency (6/13 positions), the lowest test--retest reliability (ICC $= 0.250$), and is susceptible to electromyographic and microsaccade artifacts that standard preprocessing cannot fully remove. Gamma associations that replicate in artifact-controlled modalities (MEG, intracranial EEG) would be considerably more convincing than the scalp EEG results reported here.

\subsubsection{Five-Year Test--Retest Reliability of Spectral Differentiation (Dortmund, $N = 208$)}
\label{sec:testretest}

Intraclass correlation coefficients (ICC; two-way mixed, absolute agreement) were computed between session-1 and session-2 eyes-closed enrichment profiles in the longitudinal Dortmund subsample ($N = 208$; mean inter-session interval: approximately 5 years). Per-band median ICCs are reported in Table~\ref{tab:icc}.

Beta-low enrichment was the most temporally stable band (median ICC $= 0.604$), followed by beta-high ($0.507$), alpha ($0.454$), theta ($0.382$), and gamma ($0.250$, the lowest); the cross-band median was $0.421$. At the individual-feature level, beta-low \textit{ushape} achieved ICC $= 0.746$, the highest five-year stability in the entire feature set.

Group-level profile stability was evaluated by correlating the 12-position enrichment vectors between sessions for each band. All five bands exceeded $r = 0.96$: beta-low $r = 0.988$, beta-high $r = 0.987$, theta $r = 0.983$, alpha $r = 0.977$, gamma $r = 0.964$. The pattern of eyes-closed to eyes-open change (EC$\to$EO delta vectors) was itself stable across the five-year interval ($r = 0.875$--$0.990$). Age did not predict five-year change in any enrichment feature (0 FDR across 90 tests).

\begin{table}[htbp]
\centering
\caption{Five-year test--retest ICCs for $\phisym$-lattice enrichment, Dortmund longitudinal ($N = 208$). Earlier extraction pipelines with restrictive peak caps produced negative ICCs ($-0.25$ to $-0.36$); the corrected pipeline yields strongly positive values.}
\label{tab:icc}
\begin{tabular}{lrr}
\toprule
\textbf{Band} & \textbf{Median ICC} & \textbf{Most Stable Feature (ICC)} \\
\midrule
Beta-low  & $0.604$ & beta\_low\_ushape ($0.746$) \\
Beta-high & $0.507$ & -- \\
Alpha     & $0.454$ & -- \\
Theta     & $0.382$ & -- \\
Gamma     & $0.250$ & -- \\
\midrule
\textbf{Overall} & $\mathbf{0.421}$ & \\
\bottomrule
\end{tabular}
\end{table}

\begin{figure}[htbp]
\centering
\includegraphics[width=\textwidth]{fig8_reliability.png}
\caption{\textbf{Five-year test--retest reliability.} (A)~Median intraclass correlation coefficient (ICC; two-way mixed, absolute agreement) per band across $\sim$5 years (Dortmund longitudinal subsample, $N = 208$ subjects, eyes-closed). Beta-low is the most stable band (median ICC $= 0.604$); the single most stable feature is beta-low ushape (ICC $= 0.746$, diamond). Gray dotted line: cross-band median ($0.421$). (B)~Group enrichment profile correlations (Pearson $r$) between session-1 and session-2 12-position enrichment vectors. All bands exceed $r = 0.96$, indicating that the population-level enrichment pattern is virtually identical 5 years apart.}
\label{fig:reliability}
\end{figure}

\subsubsection{Developmental Trajectory (HBN, $N = 927$, Ages 5--21)}
\label{sec:development}

All 90 enrichment features were correlated with age (continuous, in years) in the HBN pediatric dataset. Sixty tests survived FDR correction ($q < 0.05$ across 90 tests).

The strongest developmental feature was alpha \textit{inv\_noble$_4$} ($\rho = +0.354$, $p_\text{FDR} < 0.0001$): the upper flank of the alpha mountain broadens progressively throughout development, extending enrichment toward higher-frequency positions. This asymmetric broadening is corroborated by alpha \textit{asymmetry} ($\rho = +0.351$) and alpha \textit{ramp\_depth} ($\rho = +0.328$). Simultaneously, the beta-low center clears (noble$_3$: $\rho = -0.194$, $p_\text{FDR} < 0.0001$) and gamma ramp sharpens modestly (noble$_3$: $\rho = -0.207$).

Cross-release stability was assessed by correlating developmental effect vectors across the five independent HBN data releases (R1--R6, $N = 136$--322 per release). The mean inter-release Pearson correlation was $r = 0.787$. Alpha showed 13/13 position agreement on the sign of enrichment across release $\times$ position cells.

A caveat is warranted: IAF correlates with age substantially more strongly in HBN ($\rho = +0.362$, $p < 10^{-29}$, $N = 915$) than in the adult LEMON sample where IAF-partialing was performed. This IAF--age effect is comparable in magnitude to the top enrichment--age correlation reported above ($\rho = +0.354$ for alpha inv\_noble$_4$), indicating that a substantial fraction of the developmental enrichment signal may be mediated through IAF maturation rather than genuine within-band reorganization. IAF-partialed developmental correlations are not presented for HBN because the strong IAF--age collinearity ($\rho = 0.362$) makes partialing unstable: removing IAF variance necessarily removes most age variance and vice versa, yielding uninterpretable partial correlations. In the adult LEMON sample, IAF--age collinearity is weaker, permitting clean separation. Features that are most IAF-sensitive in the LEMON partialing analysis (beta-low ramp metrics, Section~\ref{sec:cognitive}) should be interpreted with proportionally greater caution in the developmental context than IAF-insensitive features (gamma ramp, theta boundary convergence).

\subsubsection{Aging Trajectory (Dortmund, $N = 608$, Ages 20--70)}
\label{sec:aging}

In the adult Dortmund sample spanning ages 20--70, 41 tests survived FDR correction. The direction of every major age effect was opposite to the developmental direction in HBN. During development, alpha \textit{inv\_noble$_3$} increases ($\rho = +0.301$); during adult aging it decreases ($\rho = -0.276$, $p_\text{FDR} < 0.001$)---the alpha mountain progressively narrows. During development, the beta-low interior clears; during adult aging, the beta-low attractor fills in ($\rho = +0.311$, $p_\text{FDR} < 0.001$)---the ascending ramp flattens.

\subsubsection{Inverted-U Lifespan Trajectory}
\label{sec:lifespan}

Combining the pediatric (HBN) and adult (Dortmund) datasets reveals a systematic inverted-U trajectory of spectral differentiation across the human lifespan. Of 31 enrichment features that reached FDR significance in at least two of the three datasets (HBN, Dortmund, LEMON), 28 showed opposite signs in the developmental versus aging comparisons ($28/31 = 90\%$). Under development, enrichment profiles become more differentiated---peaks concentrate more tightly at canonical positions; under aging, they progressively de-differentiate.

To estimate the vertex of this trajectory, we fit quadratic models to the combined HBN + Dortmund age range (5--70 years, $N = 1{,}465$--$1{,}535$ depending on metric). The quadratic model was decisively superior to linear for the three strongest features ($\Delta$BIC $> 45$). The vertex---the age at which spectral differentiation is maximal---fell consistently at \textbf{33--35 years}: alpha asymmetry 35.0 years (95\% bootstrap CI [33.5, 36.6], $R^2 = 0.093$), alpha inv\_noble$_4$ 35.2 years ([33.3, 37.5]), and beta-low ushape 33.2 years ([31.2, 35.2]). This is later than the ~20-year estimate suggested by the dataset boundary alone, and coincides more closely with the completion of cortical myelination (which continues into the early 30s) and peak cognitive performance in fluid reasoning tasks. However, the vertex estimate should be interpreted with caution: it falls near the splice point of the two contributing datasets (HBN ends at 21, Dortmund begins at 20), which have different recording systems, countries, and demographic compositions. Quadratic fits to Dortmund alone (ages 20--70, $N = 483$--507) do not support an inverted-U within that cohort: $\Delta$BIC favors linear over quadratic for all three features, and the Dortmund trajectory is monotonically declining throughout the adult range. The inverted-U is therefore driven entirely by the cross-cohort splice between developmental increase (HBN) and adult decline (Dortmund), not by a within-cohort turning point. The bootstrap CIs on the combined vertex quantify sampling uncertainty but do not address this between-dataset confound. Longitudinal data spanning childhood through middle age within a single cohort would be needed to confirm whether 33--35 years is the genuine vertex or an artifact of the cross-sectional splice.

\begin{figure}[htbp]
\centering
\includegraphics[width=\textwidth]{fig7_lifespan.png}
\caption{\textbf{Inverted-U lifespan trajectory of spectral differentiation.} (A)~Spearman rank correlations ($\rho$) with age for key enrichment features in the developmental (HBN, $N = 927$, ages 5--21, solid bars) and aging (Dortmund, $N = 608$, ages 20--70, hatched bars) datasets. FDR correction: Benjamini--Hochberg $q < 0.05$ across 90 tests per dataset; 60 survivors in HBN, 41 in Dortmund. Every major feature shows opposite signs: differentiation increases during development and decreases during aging. (B)~Schematic inverted-U trajectory. Quadratic fits (ordinary least squares) to the combined 5--70-year age range ($N = 1{,}465$--$1{,}535$) place the vertex at $\sim$33--35 years (95\% bootstrap CIs: alpha asymmetry [33.5, 36.6]; see text), coinciding with the completion of cortical myelination.}
\label{fig:lifespan}
\end{figure}

\subsubsection{Cognitive Correlates (LEMON, $N = 203$)}
\label{sec:cognitive}

Spearman correlations were computed between all 90 enrichment features and eight neuropsychological scores from the LEMON battery (eyes-closed condition) and corrected for multiple comparisons using the Benjamini--Hochberg false-discovery-rate (FDR) procedure at $q < 0.05$ across $720$ simultaneous tests ($90$ features $\times$ $8$ scores). Thirty-one tests survived FDR correction, spanning four cognitive domains and four frequency bands (Table~\ref{tab:cog_top}).

The strongest association was between beta-low \textit{center\_depletion} and logical reasoning (LPS): $\rho = -0.273$, $p_\text{FDR} = 0.024$. The negative sign indicates that individuals whose beta-low band is more deeply cleared in its interior---a steeper ascending ramp toward the $\phisym$-lattice upper boundary---achieve higher reasoning scores. The second-ranked association linked theta \textit{ushape} with the TAP incompatibility reaction-time score ($\rho = +0.267$, $p_\text{FDR} = 0.024$): greater theta spectral differentiation corresponds to faster executive conflict resolution. Upper-gamma enrichment at the inv\_noble$_4$ position predicted verbal fluency on the Regensburg Word Test (RWT; $\rho = +0.260$, $p_\text{FDR} = 0.024$), and alpha \textit{inv\_noble$_6$} predicted Trail-Making Test performance ($\rho = -0.253$, $p_\text{FDR} = 0.024$).

To test whether the cognitive associations survive replication across vigilance states, the same analysis was repeated on the eyes-open (EO) condition ($N = 202$; the slight difference from the EC sample of 203 reflects subjects whose EO recordings passed quality thresholds while their EC recordings did not, or vice versa). Twenty-five tests survived FDR correction under EO, encompassing the same four cognitive domains and the same three bands (beta-low, theta, gamma); the top effect was \textit{beta\_low\_mountain} $\times$ LPS ($\rho = -0.312$, $p_\text{FDR} = 0.004$). The cognitive signal is therefore not specific to the eyes-closed state.

\paragraph{Age-partialed analysis.}
The LPS correlation attenuated from $\rho = -0.273$ to $\rho_\text{partial} = -0.153$ ($p = 0.030$) after partialing out age, a 44\% reduction. The effect survives age control, indicating that approximately half the association is genuine cognitive variance and approximately half is mediated through age-related sharpening of spectral differentiation (Section~\ref{sec:development}).

\paragraph{IAF-partialed analysis.}
The age-partialing analysis leaves open whether the cognitive signal is driven by individual differences in $\fzero$ placement rather than genuine within-band organization---since subjects whose IAF deviates from the population $\fzero = 7.60$~Hz have systematically misaligned boundaries that could masquerade as low differentiation. To test this, we computed each subject's IAF as the power-weighted centroid of specparam-detected alpha-band peaks (mean $= 9.83 \pm 0.67$~Hz, range 8.27--11.61~Hz) and repeated the cognitive correlations partialing out IAF. Among 46 features with raw $|\rho| > 0.20$, IAF-partialing attenuated mean effect size by 23\% (from 0.240 to 0.185), with 31/46 surviving $p < 0.05$. By comparison, age-partialing produced 47\% attenuation, and combined age+IAF partialing produced 51\%---meaning IAF adds only 4 percentage points of attenuation beyond age alone, reflecting the substantial collinearity between IAF and age across LEMON's 20--77-year range. The attenuation was band-specific: gamma features were nearly unaffected (gamma\_inv\_noble$_4$ $\times$ LPS: $\rho = +0.287 \to +0.273$), consistent with IAF having no direct influence on gamma-band boundaries, while beta-low features showed larger reductions (center\_depletion $\times$ LPS: $-0.248 \to -0.127$), consistent with the alpha mountain's upper flank partly determining beta-low ramp position. Spectral differentiation captures within-band organizational variance that is largely independent of where an individual's alpha peak sits.

\paragraph{Personality null.}
The same 90 features were correlated with all 133 personality and emotion subscales available in the LEMON battery (11,970 tests). Zero tests survived FDR correction. Spectral differentiation predicts cognitive performance but is psychometrically silent for personality traits.

\begin{table}[htbp]
\centering
\caption{Top FDR-significant enrichment $\times$ cognition correlations, LEMON eyes-closed ($N = 203$, $q < 0.05$ FDR across 720 tests). LPS = logical reasoning; TAP = incompatibility reaction time; RWT = verbal fluency; TMT = Trail Making Test.}
\label{tab:cog_top}
\begin{threeparttable}
\begin{tabular}{llrr}
\toprule
\textbf{Test} & \textbf{Feature} & $\boldsymbol{\rho}$ & $\boldsymbol{p_\text{FDR}}$ \\
\midrule
LPS   & beta\_low\_center\_depletion  & $-0.273$ & $0.024$ \\
TAP   & theta\_ushape                  & $+0.267$ & $0.024$ \\
TAP   & theta\_boundary                & $+0.263$ & $0.024$ \\
TAP   & theta\_mountain                & $-0.260$ & $0.024$ \\
RWT   & gamma\_inv\_noble\_4           & $+0.260$ & $0.024$ \\
RWT   & gamma\_ramp\_depth             & $+0.256$ & $0.024$ \\
TMT   & alpha\_inv\_noble\_6           & $-0.253$ & $0.024$ \\
LPS   & beta\_low\_mountain            & $-0.250$ & $0.024$ \\
LPS   & beta\_low\_ushape              & $+0.249$ & $0.024$ \\
LPS   & beta\_low\_attractor           & $-0.245$ & $0.024$ \\
LPS   & gamma\_inv\_noble\_4           & $+0.243$ & $0.025$ \\
LPS   & beta\_low\_boundary            & $+0.240$ & $0.026$ \\
\bottomrule
\end{tabular}
\begin{tablenotes}
  \small
  \item Age-partialed top correlation: $\rho_\text{partial} = -0.153$, $p = 0.030$.
        IAF-partialed: 23\% mean attenuation; 31/46 features with $|\rho| > 0.20$ survive $p < 0.05$.
        EO replication: 25 FDR survivors with peak $\rho = -0.312$.
\end{tablenotes}
\end{threeparttable}
\end{table}

\begin{figure}[htbp]
\centering
\includegraphics[width=0.65\textwidth]{fig6_cognitive.png}
\caption{\textbf{Top 12 FDR-significant cognition correlations.} Spearman rank correlations ($\rho$) between enrichment features and cognitive test scores (LEMON eyes-closed, $N = 203$ subjects, Benjamini--Hochberg FDR $q < 0.05$ across 720 simultaneous tests). Green bars: positive correlations (higher enrichment = better performance); red bars: negative correlations (steeper ramp = better performance). Error bars are not shown as these are rank correlations; exact $p_\text{FDR}$ values are reported in Table~\ref{tab:cog_top}. Four cognitive domains (LPS reasoning, TAP executive control, RWT verbal fluency, TMT trail-making) and four frequency bands (beta-low, theta, gamma, alpha) are represented.}
\label{fig:cognitive}
\end{figure}

\subsubsection{Psychopathology (HBN, $N = 906$)}
\label{sec:psychopathology}

Enrichment features were correlated with dimensional psychopathology scores (externalizing, internalizing, p-factor, and attention subscales) derived from the Child Behavior Checklist (CBCL).

\textbf{Externalizing.} Eighteen features survived FDR correction, all in the direction of \textit{flatter} enrichment profiles: gamma \textit{inv\_noble$_4$} ($\rho = -0.159$) and beta-low \textit{inv\_noble$_1$} ($\rho = +0.127$) represent reduced spectral differentiation across multiple bands. Higher externalizing scores are associated with less concentrated, less differentiated spectral architecture.

\textbf{Internalizing.} Seven features survived FDR correction. Critically, the direction was dissociable from externalizing: internalizing was associated with \textit{more} upper-alpha enrichment (alpha\_inv\_noble$_4$: $\rho = +0.126$), the opposite of the externalizing pattern in that band. Steiger's $z$-tests for dependent correlations confirmed that the externalizing and internalizing associations were significantly different in direction for all five features tested (all $z > 2.9$, all $p < 0.004$; externalizing--internalizing correlation $r = -0.22$, $N = 906$). The dissociation is formally significant, not merely descriptive.

\textbf{General and attentional factors.} The p-factor and attention subscales each yielded zero FDR survivors, confirming that enrichment-psychopathology associations are specific to the externalizing--internalizing dimension and are not driven by general symptom severity.

\subsubsection{Informative Nulls}
\label{sec:nulls}

Every falsification criterion was satisfied. Personality and emotion subscales: 0/11,970 FDR (LEMON). Medical and metabolic markers: 0 FDR across approximately 2,500 tests (LEMON). Handedness: 0 FDR (HBN continuous and dichotomized; Dortmund). Sex $\times$ age interaction: 0 FDR (HBN and Dortmund). State $\times$ age interaction: 0 FDR (LEMON and Dortmund).

The contrast between 31 cognitive FDR survivors and zero personality FDR survivors, across the same 90 features and the same subjects, establishes that enrichment features are selectively predictive of cognitive performance, not globally correlated with individual-differences variables.

\subsubsection{Cross-Band Coupling}
\label{sec:coupling}

All pairwise Spearman correlations among the 90 enrichment features were computed within each dataset. In HBN ($N = 927$), 93 cross-band feature pairs survived FDR correction. The dominant coupling was between alpha and beta-low: alpha \textit{noble$_1$} $\times$ beta-low \textit{ushape} reached $\rho = +0.412$ ($p_\text{FDR} < 0.0001$) in HBN and replicated at $\rho = +0.344$ ($p_\text{FDR} < 0.001$) in Dortmund. Individuals with taller, more concentrated alpha mountains simultaneously exhibit deeper, more differentiated beta-low ascending ramps. This alpha--beta-low coupling is interior-only (involving no boundary positions), ruling out boundary-artifact cross-band coupling reported previously \parencite{lacy2026frontiers}. The coupling product was individually stable within-session ($r = +0.393$, $p < 10^{-22}$). By contrast, beta-low and gamma features were essentially uncoupled ($\rho \approx 0$).

\subsubsection{Cross-Release Consistency (HBN, Five Independent Releases)}
\label{sec:crossrelease}

Enrichment profiles were computed independently for each HBN release ($N = 136$--322 per release). Alpha showed 13/13 position-$\times$-release agreements on enrichment sign, and the inv\_noble$_1$ position had a standard deviation of only 1.0 percentage points across releases---the most reproducible cross-sample finding. The cross-release developmental effect vector had a mean inter-release Pearson correlation of $r = 0.787$, indicating that the developmental trajectory is stable across independent acquisition batches.

\subsubsection{State Sensitivity and State Invariance}
\label{sec:state}

Theta was the most state-sensitive band: the boundary\_hi position enrichment was $+128\%$ (EC) versus $+48\%$ (EO) in Dortmund, a $\Delta$ of 80 pp. This theta state-sensitivity reflects $\fzero$ convergence: under eyes-closed conditions, theta peaks concentrate strongly toward 7.60 Hz. Beta-low ramp contrast attenuates under EO (shape preserved, magnitude reduced). Alpha mountain sharpens under EC (attractor: $+41\%$ EC versus $+36\%$ EO).

Gamma enrichment was state-invariant: the maximum EC--EO difference at any gamma position was 14 pp. State sensitivity was age-independent: zero FDR survivors when EC--EO delta magnitude was regressed on age in either LEMON or Dortmund (90 tests each). The selective responsiveness of theta and the invariance of gamma are properties of the frequency architecture itself, not epiphenomena of age-related vigilance differences.

%------------------------------------------------------------------------------
\subsection{Part IV: Method-Independence Check}
\label{sec:irasa}
%------------------------------------------------------------------------------

All preceding results use FOOOF/specparam for aperiodic-periodic separation. To test whether the enrichment landscape depends on this specific parametric approach, we replicated the full extraction and analysis pipeline using IRASA (Irregular Resampling Auto-Spectral Analysis), a non-parametric method that isolates the aperiodic component through resampling at irrational factors rather than fitting a parametric $1/f$ model \parencite{wen2016irasa}. Band-adaptive resampling factors were selected following \textcite{gerster2022irasa} to avoid highpass edge and noise floor contamination; all downstream analyses (band assignment, Voronoi enrichment, power filtering, per-subject metrics) were identical.

Three results establish selective method-independence. First, the beta-high descending ramp achieved 12/13 mean-level sign agreement with FOOOF and 13/13 cross-dataset consistency under IRASA (versus 5/13 under FOOOF), with a steeper gradient (192 versus 113 percentage points). Beta-high is more robust under IRASA than under FOOOF. Second, the core cognitive association replicated: beta-low ramp $\times$ LPS reasoning reached $\rho = -0.294$ under IRASA ($p_\text{FDR} < 0.02$), comparable to FOOOF's $\rho = -0.278$ ($p_\text{FDR} < 0.02$). IRASA produced 8 cognitive FDR survivors overall (versus 31 under FOOOF). A peak-yield subsample test established that this gap is largely attributable to IRASA's reduced per-subject peak yield rather than to fundamentally different enrichment patterns: IRASA yields only 46\% as many theta peaks and 63--65\% as many alpha and beta-low peaks per subject as FOOOF (though beta-high and gamma are closer to parity at 71--82\%). When FOOOF peaks were randomly subsampled to match IRASA's per-band yield ratios ($N = 200$ iterations), the mean FDR survivor count dropped by approximately 40\%, substantially narrowing the gap; 30\% of subsample iterations produced $\leq$2 FDR survivors, comparable to the IRASA result. The FOOOF--IRASA cognitive gap is thus primarily a statistical power artefact of lower peak yield in the bands that carry the strongest cognitive signal, rather than evidence that FOOOF's enrichment patterns are parametric artefacts. Third, developmental trajectories were equally robust: both methods yielded 68/90 FDR-significant age correlations in HBN, dominated by the same alpha features, with IRASA producing larger top effect sizes ($\rho$ up to $0.417$ versus $0.362$).

Not all bands replicated. Theta enrichment profiles were substantially inverted between methods (5/13 sign agreement)---an expected failure mode when IRASA's evaluated frequency range extends below the highpass filter edge \parencite{gerster2022irasa}. Gamma remained inconsistent under both methods. The enrichment landscape is selectively method-independent: the strongest structural, cognitive, and developmental findings survive fundamentally different aperiodic decomposition approaches; low-frequency and high-frequency band profiles remain method-sensitive.
